\section{The two-level reduction and the variational identity}\label{sec:exchangeable}
We include the structural reduction of Lee and Kim \cite[Proposition 2 and Remark 1]{LeeKim}. Let $e_k$ denote the $k$th elementary symmetric polynomial, with $e_j=0$ for indices outside the degree range and $e_0=1$.

\begin{proposition}[Two-level reduction; Lee--Kim]\label{prop:two_level}
Fix $1\le k<N$ and $\rho\ge0$. Every maximizer of
\[
 y\longmapsto e_k(e^{y_1},\ldots,e^{y_N})
 \quad\text{on}\quad
 \left\{y\in\R^N:\sum_i y_i=0,\ \sum_i y_i^2=\rho^2\right\}
\]
has at most two distinct coordinate values.
\end{proposition}
\begin{proof}
Compactness gives a maximizer. The claims for $\rho=0$ and $N=2$ are immediate. Suppose that a maximizer has three distinct coordinates $x_1,x_2,x_3$. Fix all other coordinates $R$ and set $s=x_1+x_2+x_3$, $q=x_1^2+x_2^2+x_3^2$. On the circle defined by these two constraints, the part of the objective depending on the triple is
\[
 A\sum_{i=1}^3e^{x_i}+B\sum_{i=1}^3e^{-x_i},
 \qquad A=e_{k-1}(e^R),\quad B=e^s e_{k-2}(e^R).
\]
Indeed, the terms involving none or all three coordinates are constant, and $e^{x_1+x_2}+e^{x_1+x_3}+e^{x_2+x_3}=e^s\sum_i e^{-x_i}$. The degree range $1\le k<N$ gives $A,B\ge0$ and $A+B>0$.

Put $\phi(x)=Ae^x+Be^{-x}$. At the constrained maximum, Lagrange multipliers give $\phi'(x_i)=\alpha+\beta x_i$. Define
\[
 h(x)=\phi'(x)-\alpha-\beta x,\qquad p(x)=\prod_{i=1}^3(x-x_i),\qquad
 \eta_i=\frac1{p'(x_i)}.
\]
The identities $\sum_i\eta_i=\sum_i x_i\eta_i=0$ show that $\eta$ is tangent to the constraint circle. The second-order necessary condition is
\begin{equation}\label{eq:twolevel_second}
 \sum_{i=1}^3\frac{h'(x_i)}{p'(x_i)^2}\le0.
\end{equation}
Consider the Hermite interpolation polynomial
\[
 H(x)=\sum_{i=1}^3h'(x_i)(x-x_i)
 \left[\frac{p(x)}{(x-x_i)p'(x_i)}\right]^2.
\]
It satisfies $H(x_i)=h(x_i)=0$ and $H'(x_i)=h'(x_i)$. Repeated Rolle's theorem applied to these three double zeros gives a point $\xi$ between the smallest and largest $x_i$ such that
\[
 5!\sum_{i=1}^3\frac{h'(x_i)}{p'(x_i)^2}
 =H^{(5)}(\xi)=h^{(5)}(\xi)=Ae^\xi+Be^{-\xi}>0.
\]
This contradicts \eqref{eq:twolevel_second}.
\end{proof}

\begin{proposition}[Exact variational identity; Lee--Kim]\label{prop:variational}
The constant in \eqref{eq:exchange-intro} satisfies
\begin{equation}\label{eq:LKidentity}
 \boxed{\quad C_N^\star=
 \max_{1\le k,m<N}\frac{8N}{m(N-m)}
 \sup_{t\ne0}\frac{\psi_{N,k,m}(t)}{t^2}.\quad}
\end{equation}
\end{proposition}
\begin{proof}
For $a\perp\one$, average $\exp(\lambda\langle a,x\rangle)$ over coordinate permutations and denote the resulting convex function on $[-1,1]^N$ by $\Phi(x)$. Exchangeability gives $\E e^{\lambda\langle a,X\rangle}=\E\Phi(X)$. Convexity bounds this by the largest value at a cube vertex. A vertex with $k$ positive coordinates, uniformly permuted, is the $k$th Hamming slice. If $S_k$ is a uniform $k$-subset and $y_i=2\lambda a_i$, its log-MGF is
\[
 F_{N,k}(y)=\log\E\exp\Bigl(\sum_{i\in S_k}y_i\Bigr)
 =\log\frac{e_k(e^{y_1},\ldots,e^{y_N})}{\binom Nk}.
\]
For $\lambda a\ne0$, Proposition \ref{prop:two_level} bounds this objective by that of a vector of the same norm with values $(N-m)t/N$ on $m$ coordinates and $-mt/N$ on the others. Its slice sum and squared norm are
\[
 \sum_{i\in S_k}y_i=t(H_{N,k,m}-km/N),\qquad
 \|y\|_2^2=\frac{m(N-m)}N t^2.
\]
Since $\|y\|_2^2=4\lambda^2\|a\|_2^2$, the normalized log-MGF is at most the right side of \eqref{eq:LKidentity}. Constant slices and $\lambda a=0$ contribute zero.

Conversely, for any $k,m,t$, take $X$ uniform on the $k$th Hamming slice and let $a_i=(N-m)t/(2N)$ on a fixed $m$-subset and $a_i=-mt/(2N)$ otherwise. Then
\[
 \langle a,X\rangle=t(H_{N,k,m}-km/N),\qquad
 \|a\|_2^2=\frac{m(N-m)}{4N}t^2.
\]
The defining inequality at $\lambda=1$ gives the reverse bound after taking the supremum over $t\ne0$ and the maximum over $k,m$.
\end{proof}

Applying Corollary \ref{cor:mgf} in \eqref{eq:LKidentity} gives
\begin{equation}\label{eq:globalC}
 C_N^\star\le\frac{N}{N-1}
 \max_{k,m}\frac{\min\{k(N-k),m(N-m)\}}{m(N-m)}
 =\frac{N}{N-1}.
\end{equation}
Theorem \ref{odd:thm:odd} supplies the exact odd value. Together with the sharp examples below, these results prove Theorem \ref{thm:exchange}. Substituting the projected, zero-padded weights into \eqref{eq:exchange-intro} proves Corollary \ref{cor:tail}; its tail bound follows by minimizing the Chernoff exponent.

\section{Parity, maximal variance, and explicit sharp contrasts}\label{sec:parity}
\begin{proposition}[The variance constant and the parity gap]\label{prop:variance_parity}
For $N\ge2$, define the variance constant over nonzero contrasts by
\[
 D_N=\sup_{\substack{X\in[-1,1]^N\ \mathrm{exchangeable}\\a\ne0,\ a\perp\one}}
 \frac{\Var\langle a,X\rangle}{\|a\|_2^2}.
\]
Then
\begin{equation}\label{eq:variance_constant}
 D_N=\frac{4\lfloor N^2/4\rfloor}{N(N-1)}
 =\begin{cases}N/(N-1),&N\text{ even},\\(N+1)/N,&N\text{ odd}.
 \end{cases}
\end{equation}
For even $N$, $C_N^\star=D_N$. For odd $N$, $D_N<C_N^\star<N/(N-1)$, and, with $x=1/N$,
\begin{align}
 C_N^\star-D_N
 &=\frac{2\int_0^x t\operatorname{artanh}t\,dt}
 {(1-x)\operatorname{artanh}x}>0,\label{eq:variance_gap_integral}\\
 &=\frac{2}{3N^2}+\frac{2}{3N^3}+\frac{26}{45N^4}+O(N^{-5}).\label{eq:variance_gap_expansion}
\end{align}
\end{proposition}
\begin{proof}
Exchangeability and $a\perp\one$ give
\[
 \frac{\Var\langle a,X\rangle}{\|a\|_2^2}
 =\E(X_1^2-X_1X_2)
 =\frac1{N-1}\E\sum_{i=1}^N(X_i-\bar X_N)^2.
\]
The sum $\sum_i(x_i-\bar x)^2$ is a convex function on the cube. At a vertex with $k$ positive coordinates it equals $4k(N-k)/N$, whose maximum is $4\lfloor N^2/4\rfloor/N$. A uniformly permuted nearest-balanced vertex attains the maximum for every nonzero $a\perp\one$. This proves \eqref{eq:variance_constant}.

For even $N$, \eqref{eq:globalC} and the variance lower bound prove equality. For odd $N$, use the value established in Theorem \ref{odd:thm:odd}:
\[
 C_N^\star-D_N
 =\frac{x-(1-x^2)\operatorname{artanh}x}
 {(1-x)\operatorname{artanh}x}.
\]
The numerator has derivative $2x\operatorname{artanh}x$ and vanishes at zero, proving \eqref{eq:variance_gap_integral}. Expanding $\operatorname{artanh}x$ gives \eqref{eq:variance_gap_expansion}; $\operatorname{artanh}x>x$ gives $C_N^\star<N/(N-1)$.
\end{proof}

\begin{proposition}[A sharp single-coordinate contrast]\label{prop:explicit_extremizer}
Let $X$ be uniform on the Hamming slice with $k=\lfloor N/2\rfloor$ positive coordinates, and put $a=e_1-N^{-1}\one$. Then
\[
 \|a\|_2^2=\frac{N-1}{N},\qquad
 \langle a,X\rangle=X_1-\bar X_N=2(B-p),
 \quad B\sim\Bern(p),\quad p=\frac{k}{N}.
\]
For odd $N$, the optimal bound \eqref{eq:exchange-intro} is attained at the nonzero argument $\lambda_\star=\log((N+1)/(N-1))$. For even $N$, its sharp coefficient is attained in the limit $\lambda\to0$.
\end{proposition}
\begin{proof}
For odd $N$, write $\alpha=\log((N+1)/(N-1))=\log((1-p)/p)$. The identity $1-p+pe^{2\alpha}=e^\alpha$ gives
\[
 \log\E e^{\alpha\langle a,X\rangle}
 =(1-2p)\alpha=\frac\alpha N
 =\frac{\alpha^2 L_N}{2}\frac{N-1}{N}.
\]
For even $N$, $2(B-p)$ is a symmetric sign, so
\[
 \lim_{\lambda\to0}\frac{2\log\E e^{\lambda\langle a,X\rangle}}
 {\lambda^2\|a\|_2^2}
 =\lim_{\lambda\to0}\frac{2N\log\cosh\lambda}{(N-1)\lambda^2}
 =\frac{N}{N-1}.
\]
\end{proof}

The Bernoulli optimal proxy, used below in both normalizations, is
\begin{equation}\label{eq:bernproxy}
 \sigma^2_{\mathrm{opt}}(\Bern(p))=
 \begin{cases}\displaystyle\frac{1-2p}{2\log((1-p)/p)},&p\ne1/2,\\[5pt]1/4,&p=1/2.
 \end{cases}
\end{equation}
A calculus proof appears in Appendix \ref{app:bernoulli}; see also \cite{ArbelMarchalNguyen}. The sharp examples prove the odd lower bound
\begin{equation}\label{eq:oddL}
 C_N^\star\ge L_N=\frac{2}{(N-1)\log((N+1)/(N-1))},
\end{equation}
while \eqref{eq:globalC} gives the upper bound for all $N$. Since $L_N=1+N^{-1}+O(N^{-2})$, the first-order limit \eqref{eq:asymptotic} follows. The odd upper bound and the complete hypergeometric equality classification are proved in Section \ref{sec:odd_uniform}.
